%; whizzy paragraph -pdf xpdf -latex ./whizzypdfptex.sh
%; whizzy-paragraph "^\\\\begin{frame}\\|\\\\emtext"
% latex beamer presentation.
% platex, latex-beamer でコンパイルすることを想定。

%     Tokyo Debian Meeting resources
%     Copyright (C) 2012 Junichi Uekawa

%     This program is free software; you can redistribute it and/or modify
%     it under the terms of the GNU General Public License as published by
%     the Free Software Foundation; either version 2 of the License, or
%     (at your option) any later version.

%     This program is distributed in the hope that it will be useful,
%     but WITHOUT ANY WARRANTY; without even the implied warreanty of
%     MERCHANTABILITY or FITNESS FOR A PARTICULAR PURPOSE.  See the
%     GNU General Public License for more details.
%     You should have received a copy of the GNU General Public License
%     along with this program; if not, write to the Free Software
%     Foundation, Inc., 51 Franklin St, Fifth Floor, Boston, MA  02110-1301 USA

\documentclass[cjk,dvipdfmx,12pt]{beamer}
\usetheme{Tokyo}
\usepackage{monthlypresentation}

%  preview (shell-command (concat "evince " (replace-regexp-in-string "tex$" "pdf"(buffer-file-name)) "&"))
%  presentation (shell-command (concat "xpdf -fullscreen " (replace-regexp-in-string "tex$" "pdf"(buffer-file-name)) "&"))
%  presentation (shell-command (concat "evince " (replace-regexp-in-string "tex$" "pdf"(buffer-file-name)) "&"))

%http://www.naney.org/diki/dk/hyperref.html
%日本語EUC系環境の時
\AtBeginDvi{\special{pdf:tounicode EUC-UCS2}}
%シフトJIS系環境の時
%\AtBeginDvi{\special{pdf:tounicode 90ms-RKSJ-UCS2}}

\newenvironment{commandlinesmall}%
{\VerbatimEnvironment
  \begin{Sbox}\begin{minipage}{1.0\hsize}\begin{fontsize}{8}{8} \begin{BVerbatim}}%
{\end{BVerbatim}\end{fontsize}\end{minipage}\end{Sbox}
  \setlength{\fboxsep}{8pt}
% start on a new paragraph

\vspace{6pt}% skip before
\fcolorbox{dancerdarkblue}{dancerlightblue}{\TheSbox}

\vspace{6pt}% skip after
}
%end of commandlinesmall

\title{東京エリアDebian勉強会}
\subtitle{2020年3月度} 
\author{上川純一}
\date{2020年3月21日}
\logo{\includegraphics[width=8cm]{image200607/openlogo-light.eps}}

\begin{document}

\begin{frame}
\titlepage{}
\end{frame}

\begin{frame}{会場の注意}
 \begin{itemize}
  \item 発言していないときはミュートしてください
  \item 質問などはIRCで
  \item 資料は多分手元でもみたほうがよい
 \end{itemize}
\end{frame}


\section{事前課題}

\emtext{事前課題}

{\footnotesize
 \begin{prework}{ username }
  \begin{enumerate}
  \item xxx
  \end{enumerate}
\end{prework}

}

\section{ビデオカンファレンス}
\emtext{ビデオカンファレンス}

\begin{frame}{Why}
良い機会だしオンラインで会議するには何が必要か語らないか
\end{frame}

\begin{frame}{形態}
 \begin{itemize}
  \item 単一拠点の講演
  \item 複数拠点の講演
  \item 複数拠点での議論
 \end{itemize}
\end{frame}

\begin{frame}{RTMP}
  \begin{itemize}
  \item 古の配信プロトコル
  \end{itemize}
\end{frame}

\begin{frame}{nginx rtmp module}
 \begin{itemize}
  \item サーバにつなげてください。
 \end{itemize}
\end{frame}

\begin{frame}{OBS}
 \begin{itemize}
  \item 超便利なツール、配信に活用されている。
  \item ソースを複数ミックス、RTMPサーバに送りつけることができる
 \end{itemize}
\end{frame}

\begin{frame}{vlc}
 \begin{itemize}
  \item 便利なビデオ再生ツール、rtmp配信も見ることができる
 \end{itemize}
\end{frame}

\begin{frame}{ffmpeg}
 \begin{itemize}
  \item 便利なビデオエンコーディングツール、rtmpに投げることができる
 \end{itemize}
\end{frame}

\end{document}

;;; Local Variables: ***
;;; outline-regexp: "\\([ 	]*\\\\\\(documentstyle\\|documentclass\\|emtext\\|section\\|begin{frame}\\)\\*?[ 	]*[[{]\\|[]+\\)" ***
;;; End: ***
